%%=============================================================================
%% CONCLUSIONES Y RECOMENDACIONES
%%=============================================================================
\chapter{Conclusiones y recomendaciones}
\label{cap:conclusiones}

\section{Conclusiones}

Esta tesis evaluó QPP-RNG como fuente de entropía para criptografía ligera en dispositivos IoT críticos, mediante el diseño, la implementación en Rust y la evaluación experimental de una librería que integra el generador con ASCON-AEAD128. Las conclusiones principales son las siguientes.

\begin{enumerate}
        \item \textbf{La implementación de referencia reproduce la min-entropía IID publicada, pero no el comportamiento estadístico de la salida cruda.} Su min-entropía IID es de 7.85 bits por byte y SP 800-90B (vía no IID) le asigna entre 7.2 y 7.4 bits por byte. Sin embargo, su salida cruda no supera las pruebas SP 800-22 y ENT por un sesgo estructural en el bit más significativo, consecuencia de extraer $n_p \bmod 256$ cuando $n_p$ es geométrico con media $5! = 120$, que la evaluación publicada no reporta. Ese sesgo se predice con exactitud ($0.48654$ frente a $0.48621$ medido) y obliga a acondicionar la salida antes de usarla. Acondicionada con SHA-256, supera NIST SP 800-22, ENT y las pruebas IID de SP 800-90B a la tasa esperada por azar.

    \item \textbf{En microcontroladores, el \textit{jitter} de ejecución de QPP-RNG no aporta entropía.} En un Arduino Mega 2560, un Arduino Nano, un nRF52840 y una ESP32, la salida de la referencia es idéntica bit a bit tras cada reinicio, y aun así supera las pruebas estadísticas. En procesadores sin sistema operativo ni otros maestros de bus, el tiempo de la caminata es una función determinística de la semilla, incluso cuando, como en la ESP32, el código se ejecuta desde Flash a través de una caché.

    \item \textbf{En anfitriones rápidos, la entropía existe, pero es menor que la contabilizada.} En una Raspberry Pi 4 y en un computador con Windows, la referencia no es determinística, pero cada ciclo aporta entre 0.375 y 0.407 bits de min-entropía fresca según NIST, unos 2 bits por byte de salida frente a los 8 que contabiliza el diseño. Ninguna prueba aplicada a la salida detecta este déficit.

    \item \textbf{Una fuente física independiente resuelve el problema en microcontroladores.} El perro guardián de los ATmega, medido con el Timer1, aporta 1.62 y 2.24 bits por muestra en los dos ejemplares evaluados, y los RNG del nRF52840 y de la ESP32 aportan 7.24 y 7.15 bits por byte, todos estimados con un millón de muestras. En la ESP32, el RNG requiere mantener activa la fuente de ruido del ADC, y su reloj lento, con 2.57 bits por muestra, ofrece una alternativa del mismo tipo que la de AVR. La variante \texttt{qpp-rng-iot}, que toma su entropía de esas fuentes, aplica las pruebas de salud continuas de SP 800-90B y contabiliza explícitamente la entropía, produce salidas distintas tras cada reinicio y presenta resultados estadísticos globalmente favorables. En las cuatro placas superó la batería completa sobre 1\,000\,000 de bytes, con entre 7.26 y 7.36 bits por byte en \texttt{ea\_non\_iid}, todas las pruebas IID de \texttt{ea\_iid} superadas y entre el 99.17\,\% y el 99.51\,\% de las combinaciones de SP 800-22 superadas, por encima de lo esperado por azar.

    \item \textbf{La variante IoT es viable en un microcontrolador de 8 bits.} El cómputo por byte en el ATmega2560 se redujo $6.6\times$, de modo que el generador quedó limitado por la fuente: 7.4 B/s de salida cruda con al menos 8 bits acreditados por byte, en 6.3 a 6.5 KB de \textit{firmware}. La pila completa (fuente, QPP-RNG, SHA-256 y ASCON-AEAD128) ocupa entre 15.7 y 21.0 KB de Flash con los \textit{backends} compactos y funciona de extremo a extremo en las cuatro placas. Además, cada reinicio produjo una clave distinta, y todos los textos cifrados y etiquetas coincidieron con la implementación de referencia. En el ATmega328P, la pila usa como máximo 1\,691 de los 2\,048 bytes de RAM, y la primera clave está lista 8.7~s después del arranque; en las placas con RNG en hardware, en 0.7~s (nRF52840) y 0.2~s (ESP32).

    \item \textbf{La implementación base de ASCON-AEAD128 es correcta y está desacoplada del generador.} Supera los 1\,089 vectores KAT, una validación local de ACVP de 120 casos y 6\,840 verificaciones diferenciales y negativas contra la implementación de referencia en C, sin discrepancias.
\end{enumerate}

En respuesta a la hipótesis de la investigación, QPP-RNG \emph{en su forma original} no es una fuente de entropía viable para dispositivos IoT con microcontroladores sin sistema operativo, y en anfitriones de propósito general su contabilidad de entropía es optimista por un factor cercano a cuatro. La hipótesis se confirma solo en parte: la variante propuesta es viable en microcontroladores, pero conserva únicamente la estructura de QPP-RNG, en otras palabras, su entropía proviene de una fuente física independiente y no del \textit{jitter} de ejecución de la caminata. Una adaptación que conserva la estructura de QPP, pero toma su entropía de un oscilador físicamente independiente y la contabiliza de forma explícita, sí es viable, con la condición de que el valor acreditado provenga de una evaluación de la fuente en el dispositivo concreto. En esa adaptación, la caminata de QPP actúa como etapa de mezcla y no como fuente de entropía.

\section{Recomendaciones y trabajo futuro}

\paragraph{Para completar la validación.}
\begin{itemize}
    \item En AVR, ubicar las constantes en la memoria Flash en lugar de copiarlas a la RAM, empezando por la tabla de SHA-256, para ampliar el margen de unos 350 bytes que hoy le queda a una aplicación en el ATmega328P.
    \item Medir el consumo energético por byte generado y por operación de cifrado, una métrica que el Marco teórico contempla y que este trabajo no llegó a medir.
    \item Evaluar más ejemplares de cada placa y distintas condiciones de temperatura y tensión, para estimar la variación del valor acreditado entre dispositivos.
    \item Repetir la comparación entre reinicios con ciclos completos de alimentación, y no solo con la línea de reinicio, en particular en la ESP32.
    \item Medir en la ESP32 el desempeño de Ascon-AEAD128 con la caché caliente, repitiendo cada operación, para obtener cifras comparables con las demás placas.
    \item Evaluar con un millón de muestras el RNG de la ESP32 con la fuente de ruido del ADC desactivada, para cuantificar cuánto depende su entropía de esa fuente.
\end{itemize}

\paragraph{Para completar el diseño.}
\begin{itemize}
    \item Hacer que el acondicionador acepte una fuente falible, de modo que \texttt{QppEntropy} retorne un error ante una falla de salud en lugar de entrar en pánico.
    \item Definir una estrategia de \textit{nonces} con unicidad garantizada entre reinicios, por ejemplo un contador persistente, y una política de resiembra tras una falla de salud o un tiempo prolongado de operación.
    \item Usar la tubería para sembrar un DRBG al arranque, en lugar de producir directamente cada clave y \textit{nonce}.
    \item Ofrecer constantes de valor acreditado por placa, para que quien integra la librería no tenga que introducir valores sin justificación.
    \item Implementar \texttt{rand\_core::CryptoRng} solamente cuando todas las validaciones anteriores se hayan completado.
\end{itemize}

\paragraph{Para la referencia en anfitriones.}
\begin{itemize}
    \item Aumentar el sobremuestreo a unos 22 ciclos por byte, o aplicar la contabilidad explícita de la variante IoT al temporizador del anfitrión.
    \item Estudiar si el hecho de que el reloj de la CPU y el temporizador genérico de la Raspberry Pi provengan del mismo cristal hace parcialmente predecible la fase del flanco, lo que reduciría aún más las cifras de E4.
\end{itemize}

\paragraph{Para la metodología.}
\begin{itemize}
    \item Incorporar al arnés la comparación entre reinicios y la estimación de la entropía de la fuente condicionada al estado, como pruebas de primer nivel.
    \item Endurecer la integración con STS, es decir, resolver la ruta real del ejecutable antes de buscar sus plantillas, y reportar como no válido cualquier informe en el que una prueba haya abortado o en el que SP 800-22 no haya llegado a ejecutarse.
    \item Extender la evaluación a RISC-V, otra arquitectura común en IoT, para la cual aún no existe un temporizador específico en la librería; las variantes más recientes de la familia ESP32 (por ejemplo, ESP32-C3 y ESP32-C6) usan núcleos RISC-V y permitirían reutilizar buena parte del trabajo hecho para la ESP32.
\end{itemize}
